\documentclass{article}
\usepackage{graphicx}
\usepackage[english, russian]{babel}
\usepackage{amsmath, amssymb}
\usepackage{geometry}
\geometry{a4paper, margin=2.5cm}
\usepackage{hyperref}

\title{Прецессирующий волчок}
\date{January 2026}

\begin{document}

\maketitle

\begin{abstract}
В работе рассматривается динамика осесимметричного волчка, опирающегося остриём на идеально гладкую горизонтальную плоскость. 
В отличие от классического волчка Лагранжа, точка опоры не фиксирована, что допускает скольжение ножки. 
\end{abstract}

\section{Введение}

Волчок --- это разновидность гироскопа, осесимметричное твёрдое тело. 
При достаточно быстром вращении вокруг своей вертикальной оси волчок способен сопротивляться опрокидыванию.

Волчок Лагранжа --- это модель, в которой точка опоры жёстко закреплена на шарнире. 
В этой точке возникает сила реакции, которая удерживает ножку на месте, не давая ей скользить. 
Из-за этого центр масс волчка движется по сфере. 
Известно, что эта система является интегрируемой, а волчок способен бесконечно долго вращаться и не падать при достаточно сильной начальной закрутке. 
Подробное описание этих свойств приведено в классическом учебнике Гольдштейна [1].

Еще одним примером интегрируемой системы является волчок Эйлера --- это модель, в которой волчок находится в невесомости и вращается вокруг своего неподвижного центра масс.

 В работе [2] оцениваются необходимые условия для стабилизации изначально наклоненного, закрученного незакрепленного на плоскости волчка, с учетом трения. 

В отличие от указанных подходов, в данной работе рассматривается иная задача. 
Моделируется волчок, ножка которого опирается на идеально гладкую горизонтальную поверхность, что позволяет ножке скользить по ней при вращении в силу отсутствия сил трения в ножке, а проекции центра масс на плоскость --- оставаться на месте. 
Сила реакции опоры рассчитывается таким образом, чтобы избежать вертикального движения точки касания. 

\section{Описание системы}

Для моделирования движения волчка необходимо определить постоянные и динамические величины системы.

\subsection{Постоянные величины}

Волчок характеризуется следующими постоянными параметрами:
\begin{itemize}
\item масса $m$,
\item тензор инерции относительно системы координат, жёстко связанной с телом:
\[
\mathbf{I}_0 = \begin{pmatrix}
I_{0x} & 0 & 0 \\
0 & I_{0y} & 0 \\
0 & 0 & I_{0z}
\end{pmatrix}, 
\quad \text{причём} \quad I_{0y} = I_{0x}
\]
(волчок симметричен и вытянут относительно главной оси).
\end{itemize}




\subsection{Динамические величины}


\subsubsection{Переменные состояния}
Вектор состояния системы, эволюция которого описывается дифференциальными уравнениями, включает:
\begin{itemize}
\item вектор положения центра масс в мировой системе координат $\mathbf{r_c(t)}$,
\item импульс $\mathbf{P}(t)$,
\item момент импульса $\mathbf{L}$,
\item матрицу поворота:

\[
\mathbf{R(t)} = \begin{pmatrix}
r_{xx} & r_{yx} & r_{zx} \\
r_{xy} & r_{yy} & r_{zy} \\
r_{xz} & r_{yz} & r_{zz}
\end{pmatrix}
\]
Физический смысл матрицы R(t) заключается в том, что первый столбец R(t) показывает направление оси x твердого тела при преобразовании в мировую систему координат в момент времени t. Аналогично, второй и третий столбцы матрицы R(t) представляют собой направления осей y и z твердого тела в мировой системе координат в момент времени t.

\[
\dot{\mathbf{R}}(t) =
\begin{pmatrix}
\boldsymbol{\omega}(t) \times 
\begin{pmatrix}
r_{xx} \\ r_{xy} \\ r_{xz}
\end{pmatrix} &
\boldsymbol{\omega}(t) \times 
\begin{pmatrix}
r_{yx} \\ r_{yy} \\ r_{yz}
\end{pmatrix} &
\boldsymbol{\omega}(t) \times 
\begin{pmatrix}
r_{zx} \\ r_{zy} \\ r_{zz}
\end{pmatrix}
\end{pmatrix}
\]
где ${\omega}(t)$ --- вектор угловой скорости.

Чтобы упростить выражение, рассмотрим два вектора 
\[
a =  \begin{pmatrix}
a_{x} \\ a_{y} \\ a_{z}
\end{pmatrix} &
b =  \begin{pmatrix}
b_{x} \\ b_{y} \\ b_{z}
\end{pmatrix}
\]

Тогда
\[
a \times b = \begin{pmatrix}
a_{y} \cdot b_{z} - b_{y} \cdot a_{z}\\
- a_{x} \cdot b_{z} + b_{x} \cdot a_{z}\\
a_{x} \cdot b_{y} - b_{x} \cdot a_{y}
\end{pmatrix}
\]

Введем матрицу
\[
a^* = \begin{pmatrix}
0      & -a_{z} & a_{y}\\
a_{z}  &    0   & -a_{x}\\
-a_{y} & a_{x}  &   0
\end{pmatrix}
\]
Таким образом
\[
a^*b = \begin{pmatrix}
a_{y} \cdot b_{z} - b_{y} \cdot a_{z}\\
- a_{x} \cdot b_{z} + b_{x} \cdot a_{z}\\
a_{x} \cdot b_{y} - b_{x} \cdot a_{y}
\end{pmatrix} = 
a \times b
\]
$a^*$ --- это кососимметрическая матрица, соответствующая вектору $a$.

Теперь производную матрицы поворота по времени можно представить в виде
\[
\dot{\mathbf{R}}(t) = \boldsymbol{\omega}(t)^* {\mathbf{R}}(t)
\]
\end{itemize}


\subsubsection{Вычисляемые величины}
Остальные величины выражаются через переменные состояния:
\begin{itemize}
\item скорость центра масс
\[
\mathbf{v}(t) = \frac{\mathbf{P}(t)}{m}.
\]

\item угловая скорость
\[
\boldsymbol{\omega}(t) = \mathbf{I}^{-1}(t) \, \mathbf{L}(t).
\]

\item внешние силы: сила тяжести $\mathbf{F}_g(t) = m\mathbf{g}$ и сила реакции опоры $\mathbf{N(t)}$ (вычисление описано далее),

\item момент сил,
\[\boldsymbol{\tau}(t) = \mathbf{r}(t) \times \mathbf{N}(t)\]
где $\mathbf{r}(t)$ --- вектор от ножки волчка к центру масс.
\item тензор инерции в мировой системе координат
\[
\mathbf{I}(t) = \mathbf{R}(t) \, \mathbf{I}_0 \, \mathbf{R}^T(t).
\]
\end{itemize}

Система дифференциальных уравнений имеет вид:


\[
\frac{d}{dt}\begin{pmatrix}
r_c(t)\\
R(t)\\
L(t)\\
P(t)\\
\end{pmatrix} = 
\begin{pmatrix}
v(t)\\
\omega(t)^*R(t)\\
\tau(t)\\
F(t)\\
\end{pmatrix}
\]



% разобраться с r ro r``  рисунок


\subsection{Сила реакции опоры}

Сила реакции опоры действует на острие ножки волчка. 
Чтобы исключить вертикальные перемещения точки касания (подпрыгивание или проваливание), необходимо выбрать $\mathbf{N}$ так, чтобы вертикальное ускорение этой точки равнялось нулю.

Пусть $\mathbf{r}_0''$ --- вектор от центра масс к ножке в связанной системе координат. 
Введём обозначение $\boldsymbol{\rho} = \mathbf{R} \mathbf{r}_0''$. 
Тогда положение точки контакта в инерциальной системе:
\begin{equation}
\mathbf{r}_p = \boldsymbol{\rho} + \mathbf{r}_c
\end{equation}

Дифференцируя по времени:
\begin{equation}
\dot{\mathbf{r}}_p = \mathbf{v}_c + \boldsymbol{\omega} \times \boldsymbol{\rho}
\end{equation}

Второе дифференцирование даёт:
\begin{equation}
\ddot{\mathbf{r}}_p = \mathbf{a}_c + \dot{\boldsymbol{\omega}} \times \boldsymbol{\rho} + \boldsymbol{\omega} \times (\boldsymbol{\omega} \times \boldsymbol{\rho}) \label{eq:rp_ddot}
\end{equation}

На волчок действуют сила реакции $\mathbf{N} = N \hat{\mathbf{z}}$ (горизонтальные составляющие отсутствуют из-за гладкости поверхности) и сила тяжести $m\mathbf{g} = -mg\hat{\mathbf{z}}$. 
По второму закону Ньютона для центра масс:
\begin{equation}
\mathbf{a}_c = -g \hat{\mathbf{z}} + \frac{N}{m} \hat{\mathbf{z}} \label{eq:ac}
\end{equation}

Выразим производную угловой скорости. 
Используя $\mathbf{L} = \mathbf{I} \boldsymbol{\omega}$ и $\dot{\mathbf{L}} = \boldsymbol{\tau}$, получаем:
\begin{equation}
\dot{\boldsymbol{\omega}} = \dot{\mathbf{I}}^{-1} \mathbf{L} + \mathbf{I}^{-1} \boldsymbol{\tau} \label{eq:omega_dot}
\end{equation}

Учитывая, что $\mathbf{I} = \mathbf{R} \mathbf{I}_0 \mathbf{R}^T$, производная обратного тензора инерции имеет вид:
\begin{equation}
\dot{\mathbf{I}}^{-1} = \boldsymbol{\omega}^* \mathbf{I}^{-1} + \mathbf{I}^{-1} \boldsymbol{\omega}^{*T} \label{eq:invI_dot}
\end{equation}
где $\boldsymbol{\omega}^*$ --- кососимметрическая матрица, соответствующая вектору $\boldsymbol{\omega}$.

Момент силы реакции опоры относительно центра масс:
\begin{equation}
\boldsymbol{\tau} = \boldsymbol{\rho} \times \mathbf{N} = N_z (\boldsymbol{\rho} \times \hat{\mathbf{z}}) \label{eq:tau}
\end{equation}

Подставляя \eqref{eq:invI_dot} и \eqref{eq:tau} в \eqref{eq:omega_dot}, находим:
\begin{equation}
\dot{\boldsymbol{\omega}} = \bigl( \boldsymbol{\omega}^* \mathbf{I}^{-1} + \mathbf{I}^{-1} \boldsymbol{\omega}^{*T} \bigr) \mathbf{L} + N \mathbf{I}^{-1} (\boldsymbol{\rho} \times \hat{\mathbf{z}}) \label{eq:omega_dot_final}
\end{equation}


Теперь подставим \eqref{eq:ac} и \eqref{eq:omega_dot_final} в \eqref{eq:rp_ddot} и спроецируем на вертикальную ось $z$:
\begin{equation}
\ddot{r}_p = -g + \frac{N}{m} + \Biggl[ \Bigl( \bigl( \boldsymbol{\omega}^* \mathbf{I}^{-1} + \mathbf{I}^{-1} \boldsymbol{\omega}^{*T} \bigr) \mathbf{L} \Bigr) \times \boldsymbol{\rho} + N \bigl( \mathbf{I}^{-1} (\boldsymbol{\rho} \times \hat{\mathbf{z}}) \bigr) \times \boldsymbol{\rho} + \boldsymbol{\omega} \times (\boldsymbol{\omega} \times \boldsymbol{\rho}) \Biggr] \cdot z \label{eq:z_ddot}
\end{equation}

Группируя члены с $N_z$ и перенося остальные в правую часть, получаем:
\begin{equation}
N \left( \frac{1}{m} + \Bigl[ \bigl( \mathbf{I}^{-1} (\boldsymbol{\rho} \times \hat{\mathbf{z}}) \bigr) \times \boldsymbol{\rho} \Bigr] \cdot z \right) = g - \Bigl[ \Bigl( \bigl( \boldsymbol{\omega}^* \mathbf{I}^{-1} + \mathbf{I}^{-1} \boldsymbol{\omega}^{*T} \bigr) \mathbf{L} \Bigr) \times \boldsymbol{\rho} + \boldsymbol{\omega} \times (\boldsymbol{\omega} \times \boldsymbol{\rho}) \Bigr] \cdot z
\end{equation}

Отсюда выражаем силу реакции опоры:
\begin{equation}
N = \dfrac{g - \left[ \left( \bigl( \boldsymbol{\omega}^* \mathbf{I}^{-1} + \mathbf{I}^{-1} \boldsymbol{\omega}^{*T} \bigr) \mathbf{L} \right) \times \boldsymbol{\rho} + \boldsymbol{\omega} \times (\boldsymbol{\omega} \times \boldsymbol{\rho}) \right] \cdot z}{\dfrac{1}{m} + \left[ \left( \mathbf{I}^{-1} (\boldsymbol{\rho} \times \hat{\mathbf{z}}) \right) \times \boldsymbol{\rho} \right] \cdot z}
\end{equation}

Данное выражение позволяет вычислять силу реакции в каждый момент времени, обеспечивая выполнение кинематического условия отсутствия вертикального перемещения точки контакта.

\subsection{Численный метод} 

Как известно одним из свойств матрицы поворота является ортогональность $\mathbf{R}(t)\,\mathbf{R}^T(t) = \mathbf{I}$. При численном интегрировании уравнения 
\[
\dot{\mathbf{R}}(t) = \boldsymbol{\omega}(t)^* {\mathbf{R}}(t)
\]
в матрице ${\mathbf{R}}(t)$ неизбежно начинают накапливаться численные ошибки, что делает матрицу поворота не ортогональной и в последствии искажает влияние ${\mathbf{R}}(t)$ на тело.

Чтобы избежать этого можно периодически использовать метод ортогонализации Грама — Шмидта, но этот подход вычислительно затратный и сам может вносить дополнительные численные ошибки в систему.

В данной работе вместо матрицы поворота для описания ориентации используется 
кватернион $\mathbf{q} = [s,\;\mathbf{v}]$ — четырёхмерный вектор единичной длины, 
представляющий вращение. В отличие от матрицы, избыточность кватерниона минимальна 
(4 параметра вместо 9), а условие корректности сводится к одному равенству 
$s^2 + v_x^2 + v_y^2 + v_z^2 = 1$. При накоплении численных ошибок кватернион 
достаточно перенормировать:
\[
\mathbf{q} \leftarrow \frac{\mathbf{q}}{\|\mathbf{q}\|}.
\]
Матрица поворота $\mathbf{R}(t)$, необходимая для вычисления тензора инерции 
в мировой системе, вычисляется из кватерниона на каждом шаге. 
Уравнение эволюции кватерниона имеет вид
\[
\dot{\mathbf{q}}(t) = \tfrac{1}{2}\boldsymbol{\omega}(t)\mathbf{q}(t),
\]
где умножение понимается как кватернионное произведение $[0,\boldsymbol{\omega}(t)]$ и $\mathbf{q}(t)$.

Подробное обсуждение преимуществ кватернионов перед матрицами поворота 
и вывод уравнения $\dot{\mathbf{q}}(t)$ приведены в работе Бараффа [3].


\begin{verbatim}
struct quaternion
{
    double s;
    double v[3];
};
\end{verbatim}

% !уточнение системы дифференциальныхуравнений!

Теперь можно описать структуру, которая будет хранить состояние волчка:
\begin{verbatim}
struct Top
{
    /* Constant quantities */
    const double m;            /* mass */
    const double I0[3][3];     /* inertia tensor in body-space */
    const double inv_I0[3][3]; /* inverse of I0 */
    const double r0[3];        /* r`` */
    double r[3];               /* r` = R*r`` */
    

    /* State variables */
    double x[3];               /* x(t) */
    quaternion q;              /* q(t) */
    double P[3];               /* P(t) */
    double L[3];               /* L(t) */

    /* Derived quantities (auxiliary variables) */
    double inv_I[3][3];        /* inverse of tensor inertia in world-space */
    double R[3][3];            /* R(t) */
    double w[3];               /* w(t) */
    double v[3];               /* v(t) */

    /* Computed quantities */
    double F[3];               /* F(t) */
    double torque[3];          /* torque(t) */

};
\end{verbatim}


Для численного решения системы обыкновенных дифференциальных уравнений используется явный метод 
Рунге–Кутты пятого порядка, коэффициенты которого взяты из работы 
Дорманда–Принса [4]. Это та же схема, что лежит в основе 
известного вложенного метода DOPRI5(4), но без использования оценки 
погрешности и адаптивного управления шагом — в данной работе интегрирование 
ведётся с постоянным шагом $\Delta t$.








\begin{thebibliography}{99}
\bibitem{goldstein} Goldstein H., Poole C., Safko J. \href{https://www.math.toronto.edu/khesin/biblio/GoldsteinPooleSafkoClassicalMechanics.pdf}{\textit{Classical Mechanics}}. 3rd ed., Addison-Wesley, 2001. 

\bibitem{Ryspek Usubamatov}Ryspek Usubamatov \href{https://link.springer.com/book/10.1007/978-3-030-99213-2}{\textit{Theory of Gyroscopic Effects for Rotating Objects}} 

\bibitem{David Baraff}David Baraff \href{https://www.cs.cmu.edu/~baraff/sigcourse/notesd1.pdf}{\textit{An Introduction to Physically Based Modeling:
Rigid Body Simulation I—Unconstrained Rigid Body
Dynamics}} 
\bibitem{DormandPrince}
J.R. Dormand, P.J. Prince.
\href{https://doi.org/10.1016/0771-050X(80)90013-3}{\textit{A family of embedded Runge-Kutta formulae}}.
J. Comput. Appl. Math., 6(1):19--26, 1980.



\end{thebibliography}

\end{document}
